\structure{ОСНОВНАЯ~ЧАСТЬ}

В основной части работы рассмотрены теоретические и практические аспекты
верификации алгоритмов методом модельной проверки трасс с использованием TLA+ и
TLC.

В разделе~1 изложены теоретические основы: описана модель вычислений в TLA+,
каноническая форма спецификаций, а также принцип проверки трасс исполнения
путём сопоставления их с поведениями формальной модели.

В разделе~2 представлена методика верификации. Формализована задача валидации
трасс. Описана реализация ограниченной спецификации (\texttt{TraceSpec}),
позволяющей свести задачу к ограниченной модельной проверке в TLC. Также
разработана библиотека инструментирования на языке C для автоматической
генерации трасс исполнения в формате NDJSON.

В разделе~3 методика продемонстрирована на примере алгоритма двухфазной
транзакции. Представлены реализация менеджера ресурсов и менеджера транзакций
на языке C, спецификация алгоритма на TLA+, а также процедура
инструментирования кода. Проведённая валидация позволила обнаружить логическую
ошибку в реализации, связанную с некорректной обработкой повторных сообщений
\texttt{Prepared}, что подтверждает практическую применимость предложенного
подхода.

\section{Теоретические основы}

В основе работы лежат методы формальной спецификации и модельной проверки
распределённых систем. Формальные методы позволяют описывать систему в виде
математической модели, задающей множество допустимых состояний и переходов
между ними, а также формулировать инварианты и свойства корректности. Однако
проверка спецификации не гарантирует корректность конкретной реализации. Для
сопоставления поведения программы с её формальной моделью используется
модельная валидация исполнения (Model-Based Trace Checking), позволяющая
проверять реальные трассы исполнения относительно формальной спецификации

\subsection{Формальная спецификация TLA+}

TLA+ — язык спецификаций, основанный на теории множеств, логике первого порядка
и темпоральной логике действий (англ. TLA, temporal logic of actions).

Темпоральная логика была введена Амиром Пнуэли в 1970-х годах \cite{amir70}.
Лесли Лэмпорт увидел недостаточность этой идеи для описания систем целиком и
пришёл к мысли о необходимости использовать конечные автоматы, которым
придавался смысл формул темпоральной логики, описывающих все возможные
пути исполнения. Таким образом родилась идея темпоральной логики действий
(TLA+), которая дополнила темпоральную логику следующим \cite{lamport08}:

\begin{itemize}
    \item инвариантность при повторении состояния;
    \item темпоральное использование кванторов существования;
    \item принятие в качестве атомарных формул не только предикатов состояния,
        но и формул действий.
\end{itemize}

TLA-спецификация — темпоральная формула, часто называемая $Spec$ и являющаяся
предикатом (утверждением) о поведении. Поведение представляет собой возможный
путь исполнения системы \cite{habrias06}.

Состоянием называется присваивание значений переменных, шагом называется пара
состояний. Теперь поведение можно представить как бесконечную последовательность
состояний, а шагами поведения можно назвать пару последовательных состояний
поведения. Предикатом состояния называется функция, результат которой —
логическое значение истина или ложь — соответствует утверждению о состоянии.
Действием называется функция, имеющая смысл предиката над шагом. В этой функции
участвуют как переменные первого шага, так и второго, которые обычно
отмечаются штрихом \cite{habrias06}.

$$
Spec \triangleq Init \wedge \Box [Next]_{(v_1, v_2, ..., v_n)}
$$

Здесь $Init$ — предикат состояния, $Next$ — действие, $v_i$ — переменные, $\Box$
— единственный в данной спецификации темпоральный оператор (истинно во всех
будущих состояниях).

TLC — это программа, которая по заданному TLA+ описанию системы и формулам
свойств перебирает состояния системы и определяет, удовлетворяет ли система
заданным свойствам.

Обычно работа с TLA+/TLC строится таким образом: описывается система в TLA+
(или уже реализованная, или та, которую собираются реализовать), формализуются
интересные свойства, запускается TLC для проверки этих свойств. Однако данная
работа ставит своей целью находить несоответствия работы программы и
спецификации, для чего и будет использоваться TLC.

\subsection{Модельная проверка трасс}

Одним из часто возникающих вопросов при использовании TLA+ является проверка
соответствия программной реализации формальной спецификации. В идеале хотелось
бы доказать, что реализация является уточнением высокоуровневой спецификации,
то есть что все возможные поведения программы допустимы спецификацией.
Однако на практике такой подход сталкивается с серьёзными проблемами
масштабируемости и применим лишь к очень небольшим программам.

Полная верификация (end-to-end verification) остаётся крайне дорогостоящей и
технически сложной задачей вне зависимости от используемых инструментов. Даже в
тех редких случаях, когда она осуществима, затраты времени и ресурсов делают её
непрактичной для большинства прикладных систем \cite{cirstea2024validating}.

Вместо этого в инженерной практике применяется более прагматичный подход —
\emph{model-based checking} или \emph{model-based trace-checking}
\cite{howard2011tracechecking}. Его идея заключается не в доказательстве того,
что \emph{все} возможные поведения реализации соответствуют спецификации, а в
проверке того, что \emph{наблюдаемые} поведения ей соответствуют.

Подход состоит из следующих этапов:

\begin{enumerate}
    \item Реализация инструментируется таким образом, чтобы фиксировать
    существенные события и изменения состояний.
    \item Во время выполнения системы формируется трасса исполнения.
    \item Полученная трасса интерпретируется как ограничение на поведение
    формальной спецификации.
    \item Средство модельной проверки (в данном случае TLC) проверяет,
    существует ли поведение спецификации, согласованное с данной трассой.
\end{enumerate}

Таким образом, задача сводится к проверке воспроизводимости наблюдаемой
последовательности переходов в рамках формальной модели.

В отличие от классической модельной проверки, здесь не требуется перебирать всё
пространство состояний системы. Проверяется лишь конкретная трасса, что
существенно снижает вычислительную сложность. При этом метод применим к
распределённым системам, где трассы могут собираться с нескольких узлов и
агрегироваться для последующего анализа.

Данный подход имеет ряд преимуществ:

\begin{itemize}
    \item низкая стоимость интеграции в существующий процесс разработки;
    \item возможность использования стандартных средств модельной проверки;
    \item обнаружение некорректных переходов даже при отсутствии
    явных сбоев выполнения;
    \item применимость к системам, написанным на различных языках
    программирования.
\end{itemize}

Следует отметить, что метод не является строгой формальной гарантией
корректности, поскольку анализируется лишь конечное множество наблюдаемых
трасс. Однако при достаточном объёме тестовых запусков вероятность обнаружения
ошибок существенно возрастает. Как показывает практика применения формальных
методов в индустрии \cite{howard2024ccf}, даже
частичная формальная проверка позволяет обнаруживать серьёзные ошибки на ранних
этапах разработки.
